%Daten zur Institution

%\input{Dozentdaten}

%\renewcommand{\fachbereich}{Fachbereich}

%\renewcommand{\dozent}{Prof. Dr. . }

%Klausurdaten

\renewcommand{\klausurgebiet}{ }

\renewcommand{\klausurtyp}{ }

\renewcommand{\klausurdatum}{ . 20}

\klausurvorspann {\fachbereich} {\klausurdatum} {\dozent} {\klausurgebiet} {\klausurtyp}

%Daten für folgende Punktetabelle

\renewcommand{\aeins}{  3  }

\renewcommand{\azwei}{  3   }

\renewcommand{\adrei}{  8  }

\renewcommand{\avier}{  7  }

\renewcommand{\afuenf}{  8  }

\renewcommand{\asechs}{  3  }

\renewcommand{\asieben}{  5  }

\renewcommand{\aacht}{  5  }

\renewcommand{\aneun}{  6  }

\renewcommand{\azehn}{  4  }

\renewcommand{\aelf}{  12  }

\renewcommand{\azwoelf}{  64  }

\renewcommand{\adreizehn}{     }

\renewcommand{\avierzehn}{    }

\renewcommand{\afuenfzehn}{    }

\renewcommand{\asechzehn}{    }

\renewcommand{\asiebzehn}{    }

\renewcommand{\aachtzehn}{    }

\renewcommand{\aneunzehn}{    }

\renewcommand{\azwanzig}{    }

\renewcommand{\aeinundzwanzig}{    }

\renewcommand{\azweiundzwanzig}{    }

\renewcommand{\adreiundzwanzig}{    }

\renewcommand{\avierundzwanzig}{    }

\renewcommand{\afuenfundzwanzig}{    }

\renewcommand{\asechsundzwanzig}{    }

\punktetabelleelf

\klausurnote

\newpage

\setcounter{section}{0}

\inputaufgabegibtloesung
{3}
{

Definiere die folgenden
\zusatzklammer {kursiv gedruckten} {} {} Begriffe.
\aufzaehlungsechs{Ein
\stichwort {Normalenfeld} {}
$N$ auf einer differenzierbaren Hyperfläche

\mavergleichskette
{\vergleichskette
{ Y
}
{ \subseteq }{ \R^n
}
{  }{
}
{    }{
}
{   }{
}
}
{}{}{.}

}{Ein
\stichwort {überdeckungskompakter} {}
topologischer Raum.

}{Der
\stichwort {Kotangentialraum} {}
in einem Punkt

\mavergleichskette
{\vergleichskette
{  P
}
{ \in }{ M
}
{  }{}
{    }{}
{   }{}
}
{}{}{}
einer
\definitionsverweis {differenzierbaren Mannigfaltigkeit}{}{}
$M$.

}{Ein
\stichwort {regulärer Punkt} {}

\mathl{Q \in L}{} zu einer
\definitionsverweis {differenzierbaren Abbildung}{}{}
\maabbdisp {\varphi} { L } { M
} {}
zwischen
\definitionsverweis {differenzierbaren Mannigfaltigkeiten}{}{}
\mathkor {} {L} {und} {M} {.}

}{Das
\stichwort {Produkt} {}
von differenzierbaren Mannigfaltigkeiten
\mathkor {} {M} {und} {N} {.}

}{Ein
\stichwort {Zusammenhang} {}
auf einem differenzierbaren Vektorbündel
\maabb {} {E} {M
} {}
über einer differenzierbaren Mannigfaltigkeit $M$.
}

}
{} {}

\inputaufgabegibtloesung
{3}
{

Formuliere die folgenden Sätze.
\aufzaehlungdrei{Der
\stichwort {Satz über die Beziehung zwischen Krümmung und Einheitsnormalenvektor einer bogenparametrisierten ebenen Kurve} {.}}{Das
\stichwort {Theorema egregium} {.}}{Der
\stichwort {Satz von Green} {}
für den Flächeninhalt.}

}
{} {}

\inputaufgabegibtloesung
{8}
{

Es sei

\mavergleichskettedisp
{\vergleichskette
{ f(x,y)
}
{ =} { ax^2+by^2+cxy
}
{ } {
}
{ } {
}
{ } {
}
}
{}{}{.}
Bestimme die
\definitionsverweis {Hauptkrümmungen}{}{}
und die
\definitionsverweis {Hauptkrümmungsrichtungen}{}{}
des
\definitionsverweis {Graphen}{}{}
zu $f$ im Nullpunkt.

}
{} {}

\inputaufgabegibtloesung
{7}
{

Beweise den Existenzsatz für den Paralleltransport auf einer Hyperfläche

\mavergleichskette
{\vergleichskette
{ Y
}
{ \subseteq }{ \R^n
}
{  }{
}
{    }{
}
{   }{
}
}
{}{}{.}

}
{} {}

\inputaufgabegibtloesung
{8}
{

Zeige, dass eine
\definitionsverweis {differenzierbare Mannigfaltigkeit}{}{}
$M$ der Dimension

\mavergleichskette
{\vergleichskette
{ n
}
{ \geq }{ 1
}
{  }{
}
{    }{
}
{   }{
}
}
{}{}{}
unendlich viele
\definitionsverweis {Diffeomorphismen}{}{}
\maabbdisp {\varphi} { M } { M
} {}
besitzt.

}
{} {}

\inputaufgabegibtloesung
{3}
{

Es sei $M$ eine
\definitionsverweis {differenzierbare Mannigfaltigkeit}{}{}
und
\maabbdisp {f} { M } {\R
} {}
eine
\definitionsverweis {stetig differenzierbare Funktion}{}{.}
Die Funktion habe im Punkt

\mavergleichskette
{\vergleichskette
{ P
}
{ \in }{ M
}
{  }{
}
{    }{
}
{   }{
}
}
{}{}{}
ein
\definitionsverweis {lokales Extremum}{}{.}
Zeige

\mavergleichskettedisp
{\vergleichskette
{ (df)_P
}
{ =} { 0
}
{ } {
}
{ } {
}
{ } {
}
}
{}{}{.}

}
{} {}

\inputaufgabegibtloesung
{5}
{

Wir betrachten die
$1$-\definitionsverweis {Differentialform}{}{}

\mavergleichskettedisp
{\vergleichskette
{ \omega
}
{ =} { -vdu+udv
}
{ } {
}
{ } {
}
{ } {
}
}
{}{}{}
auf der
$1$-\definitionsverweis {Sphäre}{}{}

\mavergleichskette
{\vergleichskette
{ S^1
}
{ \subset }{ \R^2
}
{  }{
}
{    }{
}
{   }{
}
}
{}{}{,}
wobei die Koordinaten des umgebenden $\R^2$ mit
\mathkor {} {u} {und} {v} {}
bezeichnet seien. Bestimme
\mathl{\pi^* \omega}{} unter der
\definitionsverweis {stetig differenzierbaren}{}{}
Abbildung
\maabbeledisp {\pi} {\R^2 \setminus \{0\}} { S^{1}
} {(x , y) } { { \frac{   1  }{  \sqrt{x^2 + y^2}  } } (x , y)
} {.}

}
{} {}

\inputaufgabegibtloesung
{5}
{

Beweise den Satz über die Berechnung des kanonischen Volumens auf einer riemannschen Mannigfaltigkeit $M$.

}
{} {}

\inputaufgabegibtloesung
{6 (1+1+3+1)}
{

Begründe, dass der $\R^2$ keine Struktur einer
\definitionsverweis {Mannigfaltigkeit mit Rand}{}{}
derart trägt, dass die angegebene Teilmenge $T$ der
\definitionsverweis {Rand}{}{}
ist.
\aufzaehlungvierabc{
\mavergleichskette
{\vergleichskette
{ T
}
{ = }{ {]0,1[}
}
{  }{
}
{    }{
}
{   }{
}
}
{}{}{,}
wobei das Intervall auf der $x$-Achse liegt.
}{
\mavergleichskette
{\vergleichskette
{ T
}
{ = }{ [0,1]
}
{  }{
}
{    }{
}
{   }{
}
}
{}{}{,}
wobei das Intervall auf der $x$-Achse liegt.
}{
\mavergleichskette
{\vergleichskette
{ T
}
{ = }{ \R
}
{  }{
}
{    }{
}
{   }{
}
}
{}{}{,}
wobei $\R$ die $x$-Achse sei.
}{
\mavergleichskette
{\vergleichskette
{ T
}
{ = }{ S^1
}
{  }{
}
{    }{
}
{   }{
}
}
{}{}{.}
}

}
{} {}

\inputaufgabegibtloesung
{4}
{

Es sei $\omega$ eine stetig differenzierbare
$n-1$-\definitionsverweis {Differentialform}{}{}
auf einer
\definitionsverweis {offenen Menge}{}{}

\mavergleichskette
{\vergleichskette
{ U
}
{ \subseteq }{ \R^n
}
{  }{
}
{    }{
}
{   }{
}
}
{}{}{}
und sei die
\definitionsverweis {äußere Ableitung}{}{}
gleich

\mavergleichskettedisp
{\vergleichskette
{d \omega
}
{ =} { fdx_1 \wedge \ldots \wedge dx_n
}
{ } {
}
{ } {
}
{ } {
}
}
{}{}{}
mit einer Funktion
\maabb {f} { U } {\R
} {.}
Zeige für

\mavergleichskette
{\vergleichskette
{ P
}
{ \in }{ U
}
{  }{
}
{    }{
}
{   }{
}
}
{}{}{}
die Gleichheit

\mavergleichskettedisp
{\vergleichskette
{ f(P)
}
{ =} { { \frac{   1  }{  \beta_n  } } \cdot \operatorname{lim}_{   \epsilon \rightarrow  0  } \, { \frac{   1  }{  \epsilon^n } } \int_{S^{n-1}(P, \epsilon)} \omega
}
{ } {
}
{ } {
}
{ } {
}
}
{}{}{,}
wobei $\beta_n$ das
[[Allgemeines Kugelvolumen/Mit Cavalieri-Prinzip/Beispiel|Volumen]]
der $n$-dimensionalen Einheitskugel bezeichnet.

}
{} {}

\inputaufgabegibtloesung
{12}
{

Beweise den Satz von Stokes für Mannigfaltigkeiten mit Rand.

}
{} {}
 [[Kategorie:Latexseite]]
